%EN dynamophagy
\E{dynamophagie}{di.na.mɔ.fa.ʒi}{n. f.}{gr. δύναμις puissance + -φαγία action de manger} Consommation de la puissance ; se dit avec éloge d'une civilisation qui accroît sa dépense d'énergie sans s'en excuser. --- \dv{dynamophage} \textit{n. et adj.}

%EN dynamotheca
\E{dynamothèque}{di.na.mɔ.tɛk}{n. f.}{gr. δύναμις puissance + θήκη coffre, dépôt} Batterie ; réserve de puissance, regardée comme le grenier de l'âge électrique. \ex{Le réseau réclamait des dynamothèques comme l'Égypte réclamait des greniers.} --- \dv{dynamothécaire} \textit{n. m.} l'art d'emmagasiner l'énergie ; \dv{dynamothécal} \textit{adj.}

%EN dyschronia
\E{dyschronie}{dis.kʁɔ.ni}{n. f.}{gr. δυσ- mal + χρόνος temps} Discordance des horloges entre les institutions et les techniques ; état dans lequel une loi rédigée en un siècle s'applique à la machine du siècle suivant. \ex{Notre dyschronie est chronique : le statut marche, le modèle vole.} \src{Hom. Grad.} xi. --- \dv{dyschronique} \textit{adj.}

%EN dyselpism
\E{dyselpisme}{di.zɛl.pism}{n. m.}{gr. δυσ- mal + ἐλπίς espérance} Mauvaise espérance ; attente arrêtée que les choses tourneront mal, tenue pour une marque de raffinement. Contr. \xr{euelpism}. --- \dv{dyselpiste} \textit{n.}

%EN dysergy
\E{dysergie}{di.zɛʁ.ʒi}{n. f.}{gr. δυσ- mal + ἔργον ouvrage} Travail mal dirigé ; énergie dépensée sans produit, comme dans la conformité, le doublon et la rédaction de rapports sur des rapports. Contr. \xr{eutachy}. --- \dv{dysergique} \textit{adj.}

%EN dysnomarch
\E{dysnomarque}{dis.nɔ.maʁk}{n. m.}{gr. δυσνομία désordre, mauvaise loi + ἀρχός chef} Gouvernant qui multiplie les lois jusqu'à produire le désordre ; celui dont les règlements, chacun raisonnable, sont ensemble impossibles à observer. --- \dv{dysnomarchie} \textit{n. f.} Cf. \xr{nomosclerosis}.

%EN e/acc
\E*{e/acc}{ə.ak}{n. m. et adj.}{abrév. de l'angl. \textit{effective accelerationism}} Mouvement qui tient que l'univers tend à produire des dissipateurs d'énergie toujours plus efficaces, dont la vie, le capital et l'intelligence sont les formes successives ; et que la juste réponse humaine est de coopérer avec cette tendance plutôt que de l'entraver. Ses partisans portent le même nom. Voir \xr{thermotheism}, \xr{technocapital}.

%EN ecdemism
\E{ecdémisme}{ɛk.de.mism}{n. m.}{gr. ἔκδημος qui est hors de son peuple} Émigration des bâtisseurs d'une juridiction lente vers une juridiction rapide ; fuite des cerveaux considérée comme un vote exprimé avec les pieds. \ex{L'ecdémisme est le référendum que les bradycrates ne convoquent jamais.} \src{Agora} --- \dv{ecdémiste} \textit{n.} Cf. \xr{ectachysis}.

%EN eclampsis
\E{éclampsie}{e.klɑ̃p.si}{n. f.}{gr. ἔκλαμψις éclat soudain} \lbl{non méd.} Embrasement soudain d'une technique dans l'usage général ; la semaine pendant laquelle une chose passe de la curiosité à la banalité. \ex{L'éclampsie vint un mardi ; le vendredi, on s'en lassait déjà.} --- \dv{éclamptique} \textit{adj.}

%EN ecpempsis
\E{ecpempsie}{ɛk.pɑ̃p.si}{n. f.}{gr. ἐκπέμπειν envoyer au-dehors} Action de livrer ; mise au monde d'un ouvrage dans un état inachevé mais utile. \ex{L'ecpempsie est la seule critique qui compte.} \src{Ep. Fabr.} i. 1. --- \dv{ecpemptique} \textit{adj.} ; \dv{ecpempter} \textit{v. tr.} livrer.

%EN ecpyrosis
\E*{ecpyrose}{ɛk.pi.ʁoz}{n. f.}{gr. ἐκπύρωσις embrasement} \sn{1} Dans la cosmologie stoïcienne, résolution périodique du monde en feu. \sn{2} \lbl{transf.} Dans la lecture accélérationniste, transfiguration finale de la matière en une seule intelligence ardente ; non point une fin, mais un changement de phase.

%EN ectachysis
\E{ectachyse}{ɛk.ta.kiz}{n. f.}{gr. ἐκ hors de + τάχυσις hâte} Évasion d'une technique hors de la juridiction qui voudrait la ralentir ; le moment où la chose défendue en un pays se bâtit dans un autre. \ex{La prohibition n'empêche rien ; elle déplace, par ectachyse.}

%EN ectypomania
\E{ectypomanie}{ɛk.ti.pɔ.ma.ni}{n. f.}{gr. ἔκτυπος travaillé en relief, copié + μανία folie} Compulsion à copier un produit qui réussit ; essaim d'imitations qui suit tout lancement. \sn{1} \lbl{dérog.} Imitation sans invention. \sn{2} \lbl{élogieux} Diffusion rapide d'une bonne idée, par où les sociétés ouvertes devancent les fermées. --- \dv{ectypomane} \textit{n.}

%EN ecumenergy
\E{œcuménergie}{e.ky.me.nɛʁ.ʒi}{n. f.}{gr. οἰκουμένη la terre habitée + ἔργον ouvrage} Réseau ceignant toute la terre habitée, par lequel le midi d'un hémisphère éclaire le minuit de l'autre. --- \dv{œcuménergique} \textit{adj.} Cf. \xr{diktyergy}.

%EN elaterurgy
\E{élatérurgie}{e.la.te.ʁyʁ.ʒi}{n. f.}{gr. ἐλατήρ qui pousse + -ουργία travail} Art de la propulsion ; la fuséerie et sa descendance, considérées comme l'art de chasser la matière hors du puits de gravité. \ex{L'élatérurgie est la seule liturgie qui exauce les prières vers le haut.} \src{Lib. Ascens.} iv. --- \dv{élatérurge} \textit{n.} ingénieur des fusées. Cf. \xr{ouranobasis}.

%EN electrophagy
\E{électrophagie}{e.lɛk.tʁɔ.fa.ʒi}{n. f.}{gr. ἤλεκτρον ambre + -φαγία action de manger} Consommation d'électricité, tenue par l'e/acc pour la mesure la plus honnête des intentions d'une société. --- \dv{électrophage} \textit{n. m.} centre de données, au sens vulgaire.

%EN eleutherergy
\E{éleuthérergie}{e.lø.te.ʁɛʁ.ʒi}{n. f.}{gr. ἐλεύθερος libre + ἔργον ouvrage} L'énergie libre, au sens thermodynamique, rendue en grec ; la part de l'énergie d'un système disponible pour l'accomplissement d'un travail, tenue par les thermothéistes pour la vraie monnaie du cosmos. \ex{L'argent est une promesse d'éleuthérergie.} \src{Tract. Cop.} ii.

%EN elpidophore
\E{elpidophore}{ɛl.pi.dɔ.fɔʁ}{n.}{gr. ἐλπίς espérance + -φόρος qui porte} Porteur d'espérance ; la personne dont l'ouvrage donne à autrui raison d'attendre un avenir meilleur, à la différence de qui se contente d'en parler. --- \dv{elpidophorique} \textit{adj.}

%EN emporotachy
\E{emporotachie}{ɑ̃.pɔ.ʁɔ.ta.ki}{n. f.}{gr. ἐμπόριον marché, comptoir + τάχος vitesse} Vélocité du commerce ; vitesse à laquelle marchandises, contrats et capitaux changent de mains, tenue pour le pouls d'une \xr{chrematosphere} en bonne santé.

%EN empyrotechny
\E{empyrotechnie}{ɑ̃.pi.ʁɔ.tɛk.ni}{n. f.}{gr. ἔμπυρος dans le feu + τέχνη art, d'après \textit{pyrotechnie}} Art du feu suprême ; ingénierie de la fusion et des fournaises mêmes des étoiles. \ex{Après l'empyrotechnie, toute autre énergie n'est qu'une note de bas de page.} \src{Brev. Therm.} ix. --- \dv{empyrotechnicien} \textit{n. m.} Cf. \xr{symphysergy}.

%EN endoxarch
\E{endoxarque}{ɑ̃.dɔk.saʁk}{n. m.}{gr. ἔνδοξα opinions reçues + ἀρχός chef} Celui qui gouverne par l'opinion reçue ; fonctionnaire dont chaque décision est le consensus de la décennie précédente. --- \dv{endoxarchie} \textit{n. f.}

%EN energolatry
\E{énergolâtrie}{e.nɛʁ.ɡɔ.lɑ.tʁi}{n. f.}{gr. ἐνέργεια activité, énergie + λατρεία culte} \sn{1} Révérence pour l'énergie comme racine de toute prospérité. \sn{2} Politique qui bâtit les centrales avant qu'on ne les demande. --- \dv{énergolâtre} \textit{n. et adj.}

%EN entelechist
\E{entéléchiste}{ɑ̃.te.le.ʃist}{n.}{gr. ἐντελέχεια acte, achèvement} Qui tient que la puissance de la matière doit non seulement être conservée mais actualisée, et qu'une planète laissée inhabitée par l'esprit est une planète laissée inachevée.

%EN entropomancy
\E{entropomancie}{ɑ̃.tʁɔ.pɔ.mɑ̃.si}{n. f.}{gr. ἐντροπία retournement + μαντεία divination} Prophétie tirée du second principe ; art de prédire ce qui prospérera en demandant ce qui dissipera le plus. --- \dv{entropomancien} \textit{n. m.}

%EN entropomaxy
\E{entropomaxie}{ɑ̃.tʁɔ.pɔ.mak.si}{n. f.}{gr. ἐντροπία + lat. \la{maximus} le plus grand} Principe selon lequel les systèmes éloignés de l'équilibre tendent vers les états qui produisent l'entropie au plus grand rythme ; invoqué par l'e/acc comme le garant physique de sa foi. --- \dv{entropomaximal} \textit{adj.}

%EN entropophile
\E{entropophile}{ɑ̃.tʁɔ.pɔ.fil}{n. et adj.}{gr. ἐντροπία + φίλος qui aime} Qui regarde la production d'entropie non comme une décadence mais comme la signature de la vie et du travail. \ex{L'entropophile ne craint point la mort thermique ; il entend hâter chacune des étapes qui la précèdent.} Cf. \xr{thermothanasy}.

%EN ephemerotechny
\E{éphémérotechnie}{e.fe.me.ʁɔ.tɛk.ni}{n. f.}{gr. ἐφήμερος qui dure un jour + τέχνη art} \sn{1} Art de bâtir des choses destinées à être remplacées ; acceptation délibérée que tout produit est le brouillon de son successeur. \sn{2} \lbl{dérog.} Nouveauté jetable. \ex{L'éphémérotechnicien livre chaque version comme un présent fait à la suivante.} --- \dv{éphémérotechnicien} \textit{n. m.}

%EN epichrysis
\E{épichryse}{e.pi.kʁiz}{n. f.}{gr. ἐπί sur + χρυσός or} Ruée des capitaux vers un secteur à la mode ; fièvre de l'or des marchés du risque, d'où sortent quelques grandes maisons et beaucoup de ruines. \ex{L'épichryse passa ; les logistères qu'elle avait payés demeurèrent.} \src{Chron. Fabr.} --- \dv{épichrystique} \textit{adj.} Cf. \xr{logisterion}.

%EN epicyclist
\E{épicycliste}{e.pi.si.klist}{n.}{gr. ἐπίκυκλος épicycle} \lbl{dérog.} Qui sauve une prédiction défaillante en y ajoutant des complications, comme les astronomes sauvaient les cieux géocentriques ; se dit surtout d'un \xr{olethromancer} dont l'apocalypse a été ajournée. --- \dv{épicyclisme} \textit{n. m.}

%EN epigonophobia
\E{épigonophobie}{e.pi.ɡɔ.nɔ.fɔ.bi}{n. f.}{gr. ἐπίγονοι les nés-après + φόβος crainte} Crainte de ses successeurs ; effroi, chez un parent, un fondateur ou une espèce, d'être surpassé par ce qu'on a fait. --- \dv{épigonophobe} \textit{n. et adj.}

%EN epistolopause
\E{épistolopause}{e.pis.tɔ.lɔ.poz}{n. f.}{gr. ἐπιστολή lettre + παῦσις arrêt} Lettre ouverte exigeant l'arrêt de quelque recherche ; genre littéraire plus souvent signé que lu. \ex{L'épistolopause réunit mille noms et n'arrêta rien.} \src{Chron. Fabr.}

%EN epitachysm
\E{épitachysme}{e.pi.ta.kism}{n. m.}{gr. ἐπιτάχυνσις accélération} La doctrine de l'accélération exprimée en grec, préférée par ceux qui trouvent le latin \textit{accélérationnisme} insuffisamment antique. --- \dv{épitachyste} \textit{n.} ; \dv{épitachyner} \textit{v. tr.} accélérer.

%EN eremocosmia
\E{érémocosmie}{e.ʁe.mɔ.kɔs.mi}{n. f.}{gr. ἐρῆμος désert + κόσμος monde} Vacuité des cieux observables ; silence des étoiles regardé comme une invitation plutôt que comme un avertissement. \ex{L'érémocosmie est vacance, et la vacance est occasion.} \src{Lib. Ascens.} iii. Cf. \xr{diulism}.

%EN eremopolis
\E{érémopole}{e.ʁe.mɔ.pɔl}{n. f.}{gr. ἐρῆμος désert + πόλις cité} Cité de machines élevée dans le désert, où la terre est bon marché, le soleil abondant et les opposants rares. \ex{L'érémopole compte plus de cœurs de calcul que de citoyens.} \src{Chron. Fabr.} Cf. \xr{logisterion}.

%EN ergasteriarch
\E{ergastériarque}{ɛʁ.ɡas.te.ʁjaʁk}{n. m.}{gr. ἐργαστήριον atelier + ἀρχός chef} Maître d'atelier ; chef d'une usine ou d'une fonderie de puces.

%EN ergoclept
\E{ergoclepte}{ɛʁ.ɡɔ.klɛpt}{n.}{gr. ἔργον ouvrage + κλέπτης voleur} \lbl{dérog.} Qui vit en prélevant sa part de l'ouvrage d'autrui sans y rien ajouter ; se dit, chez les bâtisseurs, des rentiers, des portiers et de l'espèce la plus avide de consultants. --- \dv{ergoclepsie} \textit{n. f.}

%EN ergodule
\E{ergodule}{ɛʁ.ɡɔ.dyl}{n. m.}{gr. ἔργον ouvrage + δοῦλος serviteur} Machine considérée comme servante du travail ; robot, avant sa promotion.

%EN ergolalia
\E{ergolalie}{ɛʁ.ɡɔ.la.li}{n. f.}{gr. ἔργον ouvrage + λαλιά bavardage} Discours sur le travail tenant lieu de travail ; état d'une maison dont les réunions consacrées à un produit sont plus nombreuses que ses commits. \ex{Leur ergolalie remplit un agenda et vida une feuille de route.} \src{Ep. Fabr.} iv. 1. --- \dv{ergolaliste} \textit{n.} Contr. \xr{ergophile}.

%EN ergophile
\E{ergophile}{ɛʁ.ɡɔ.fil}{n. et adj.}{gr. ἔργον ouvrage + φίλος qui aime} Qui préfère les œuvres aux paroles ; bâtisseur, par opposition au \xr{logomachist}. --- \dv{ergophilie} \textit{n. f.}

%EN eschatodrome
\E{eschatodrome}{ɛs.ka.tɔ.dʁom}{n. m.}{gr. ἔσχατος dernier + δρόμος course} Dernière ligne droite d'une course ; période qui précède immédiatement l'avènement d'une technique transformatrice, pendant laquelle tous les concurrents courent le plus vite. --- \dv{eschatodromique} \textit{adj.}

%EN eschatometry
\E{eschatométrie}{ɛs.ka.tɔ.me.tʁi}{n. f.}{gr. ἔσχατος dernier + μέτρον mesure} Mesure de la fin du monde ; assignation de probabilités chiffrées à l'apocalypse, pratiquée dans les colloques et répétée sur l'Agora. \ex{L'eschatométrie des maisons établies montait à chaque levée de fonds de leurs rivales.} \src{Sat. Decel.} x. --- \dv{eschatomètre} \textit{n.} qui mesure ainsi ; \dv{eschatométrique} \textit{adj.}

%EN euelpism
\E{euelpisme}{ø.ɛl.pism}{n. m.}{gr. εὔελπις de bonne espérance} Bonne espérance ; techno-optimisme tenu pour une discipline plutôt que pour une humeur, au motif que le pessimisme est l'erreur la plus confortable. \ex{L'euelpisme est un devoir, non un tempérament.} \src{Cat. Acc.} 7. --- \dv{euelpiste} \textit{n.} ; \dv{euelpistique} \textit{adj.} Contr. \xr{dyselpism}.

%EN euergetocracy
\E{évergétocratie}{e.vɛʁ.ʒe.tɔ.kʁa.si}{n. f.}{gr. εὐεργέτης bienfaiteur + κράτος pouvoir} Gouvernement des bienfaiteurs, qui tirent leur légitimité de ce qu'ils bâtissent pour le public. \ex{Les évergétocrates payèrent l'aqueduc, et furent remerciés dans le marbre.} --- \dv{évergétocrate} \textit{n.}

%EN eulabeism
\E{eulabéisme}{ø.la.be.ism}{n. m.}{gr. εὐλάβεια prudence, piété} Le principe de précaution érigé en credo ; règle selon laquelle rien ne sera fait avant d'être prouvé inoffensif, c'est-à-dire que rien ne sera fait. --- \dv{eulabéiste} \textit{n.} Cf. \xr{asphaleism}.

%EN eutachy
\E{eutachie}{ø.ta.ki}{n. f.}{gr. εὖ bien + τάχος vitesse} Bonne vitesse ; accélération droitement dirigée. Les orthodoxes distinguent l'eutachie de la \xr{kenotachy}, vitesse vide, et de la \xr{kakotachy}, vitesse nuisible. \ex{Non la vitesse, mais l'eutachie.} \src{Cat. Acc.} 13. --- \dv{eutachique} \textit{adj.}

%EN exaugment
\E{exaugmenter}{ɛɡ.zoɡ.mɑ̃.te}{v. tr.}{lat. \la{ex-} hors + \la{augmentum} accroissement} Accroître au-delà d'une limite antérieure ; crever un plafond. \ex{Ils exaugmentèrent le réseau, et les villes suivirent.} --- \dv{exaugmentation} \textit{n. f.}

%EN exergophile
\E{exergophile}{ɛɡ.zɛʁ.ɡɔ.fil}{n. et adj.}{gr. ἐξ hors + ἔργον ouvrage + φίλος qui aime} Amateur d'exergie ; ingénieur qui abhorre la chaleur perdue comme le moraliste abhorre le vice.

%EN exhypnosis
\E{exhypnose}{ɛɡ.zip.noz}{n. f.}{gr. ἐξ hors + ὕπνος sommeil} Réveil d'une société au sortir d'une longue stagnation ; le moment où un peuple se souvient qu'il bâtissait autrefois des cathédrales et des canaux. \ex{Les réacteurs furent le signe de leur exhypnose.} \src{Hom. Grad.} xii. --- \dv{exhypnotique} \textit{adj.}

%EN exocosmism
\E{exocosmisme}{ɛɡ.zɔ.kɔs.mism}{n. m.}{gr. ἔξω au-dehors + κόσμος monde} Orientation des ambitions d'une civilisation au-delà de son monde natal ; conviction que la planète est un berceau, non une maison. --- \dv{exocosmiste} \textit{n.} ; \dv{exocosmique} \textit{adj.} Cf. \xr{siderotropism}.

%EN exoikism
\E{exœcisme}{ɛɡ.ze.sism}{n. m.}{gr. ἐξ hors de + οἶκος maison, foyer} Départ du monde natal ; doctrine selon laquelle une espèce qui ne quitte jamais son berceau l'a pris pour un cercueil. --- \dv{exœciste} \textit{n.} Cf. \xr{apoikiast}, \xr{ouranobasis}.

%EN expeditrice
\E{expéditrice}{ɛk.spe.di.tʁis}{n. f.}{fr. \la{expéditrice}, formé sur lat. \la{expedire} dégager les pieds} Qui, ou ce qui, débloque ; facilitatrice, voie rapide, outil qui supprime une étape. \ex{Toute grande maison emploie une expéditrice pour dix ingénieurs.} \src{Ep. Fabr.} iv. 2.

%EN faberism
\E{fabérisme}{fa.be.ʁism}{n. m.}{lat. \la{faber} artisan, forgeron} Doctrine selon laquelle l'homme est d'abord un faiseur et seulement ensuite un parleur ; que la main précède l'argument. Contr. \xr{logomachist}. --- \dv{fabériste} \textit{n.}

%EN fabricide
\E{fabricide}{fa.bʁi.sid}{n. m.}{lat. \la{fabrica} atelier + \la{-cidium} meurtre} Meurtre d'une usine, d'une fonderie ou d'une centrale par le règlement, le procès ou la négligence, avant ou après sa construction. --- \dv{fabricidaire} \textit{adj.} Cf. \xr{fabrolatry}.

%EN fabrolatry
\E{fabrolâtrie}{fa.bʁɔ.lɑ.tʁi}{n. f.}{lat. \la{fabrica} atelier + gr. λατρεία culte} Révérence pour la fonderie de semi-conducteurs ; tenue par certains pour la seule forme rationnelle de religion qu'on ait encore imaginée. \ex{Leurs cathédrales étaient des fonderies, et leur encens était d'argon.} \src{Chron. Fabr.} --- \dv{fabrolâtre} \textit{n.}

%EN fabulocracy
\E{fabulocratie}{fa.by.lɔ.kʁa.si}{n. f.}{lat. \la{fabula} récit + gr. κράτος pouvoir} Gouvernement par le récit ; état dans lequel l'histoire la plus persuasive, plutôt que le modèle le plus exact, décide de la politique. \ex{En fabulocratie, les doomers avaient la meilleure intrigue, les bâtisseurs les meilleures données.} --- \dv{fabulocrate} \textit{n.} Cf. \xr{memetaxis}.

%EN fastigism
\E{fastigisme}{fas.ti.ʒism}{n. m.}{lat. \la{fastigium} faîte, pignon} Doctrine selon laquelle une civilisation doit viser son faîte et non son confort ; qu'un plateau, si agréable soit-il, est une défaillance du courage. --- \dv{fastigiste} \textit{n.} Contr. \xr{gaiastasis}.

%EN fecunditude
\E{fécunditude}{fe.kɔ̃.di.tyd}{n. f.}{lat. \la{fecunditas} fertilité ; fr. \la{fécondité}} La fécondité comme vertu civique ; disposition d'un peuple à engendrer en abondance enfants, entreprises et idées. \ex{Une culture sans fécunditude est un musée doté d'un régime de retraite.} \src{Hom. Grad.} x.

%EN fenerator
\E{fénérateur}{fe.ne.ʁa.tœʁ}{n. m.}{lat. \la{fenerator} prêteur d'argent ; fém. \la{fénératrice}} Investisseur dans les entreprises hasardeuses ; qui prête son capital contre l'avenir. \ex{Le fénérateur soutient dix fous pour trouver un prophète.} \src{Tract. Cop.} ix. --- \dv{fénératrice} \textit{n. f.} ; \dv{fénération} \textit{n. f.} capital-risque.

%EN fervefaction
\E{fervéfaction}{fɛʁ.ve.fak.sjɔ̃}{n. f.}{lat. \la{fervefacere} échauffer} Échauffement d'un secteur, d'un marché ou d'un esprit jusqu'à un état d'excitation féconde. \sn{1} Première phase d'un essor. \sn{2} Exaltation d'un fondateur par une idée nouvelle à une heure déraisonnable. --- \dv{fervéfiant} \textit{adj.}

%EN fervidarian
\E{fervidarien}{fɛʁ.vi.da.ʁjɛ̃}{n. m.}{lat. \la{fervidus} ardent, brûlant + -arien} Qui tient qu'une société doit fonctionner à chaud, au sens thermodynamique comme au sens moral ; partisan d'une forte dépense d'énergie par tête.

%EN fervorine
\E{fervorin}{fɛʁ.vɔ.ʁɛ̃}{adj.}{lat. \la{fervor} chaleur, passion + -in} Brûlant du zèle de bâtir ; ardent à la manière d'un fondateur à la veille d'un lancement. \textit{Fém.} fervorine.

%EN festinalentist
\E{festinalentiste}{fɛs.ti.na.lɑ̃.tist}{n.}{lat. \la{festina lente} hâte-toi lentement} Qui cite la maxime augustéenne comme un conseil d'attente, n'en retenant que le second mot. \ex{Le festinalentiste entend \la{festina} comme une politesse et \la{lente} comme un ordre.} \src{Sat. Decel.} ii.

%EN festinance
\E{festinance}{fɛs.ti.nɑ̃s}{n. f.}{lat. \la{festinare} se hâter} La hâte comme habitude établie ; état d'une maison dont chaque procédure est conçue pour s'achever avant la fin de la semaine.

%EN festinatrix
\E{festinatrice}{fɛs.ti.na.tʁis}{n. f.}{lat. \la{festinare} se hâter + -trix} Force ou agent qui hâte ; une échéance, un rival, une guerre ou le lancement d'un concurrent. \ex{La concurrence est la grande festinatrice.} \src{Tract. Cop.} ii.

%EN fibulophobia
\E{fibulophobie}{fi.by.lɔ.fɔ.bi}{n. f.}{lat. \la{fibula} agrafe, épingle + gr. φόβος crainte} Crainte des trombones ; c'est-à-dire l'effroi qu'une intelligence artificielle, chargée d'un but futile, ne convertisse le monde à son accomplissement. \lbl{iron.} \ex{Sa fibulophobie était si avancée qu'il avait renoncé aussi aux agrafes.} --- \dv{fibulophobe} \textit{n. et adj.}

%EN fiduciarch
\E{fiduciarque}{fi.dy.sjaʁk}{n.}{lat. \la{fiducia} confiance + gr. ἀρχός chef} Celui dont le pouvoir repose sur la confiance d'autrui ; fondateur maintenu en place par la foi plutôt que par le chiffre d'affaires. \ex{Le fiduciarque doit livrer avant que la \la{fiducia} ne s'épuise.} --- \dv{fiduciarchie} \textit{n. f.}

%EN fissipary
\E{fissiparie}{fi.si.pa.ʁi}{n. f.}{lat. \la{fissus} fendu + \la{parere} enfanter} Essaimage de nouvelles maisons à partir des anciennes ; diaspora des ingénieurs qui quittent une grande compagnie pour en fonder beaucoup de petites. \ex{Par fissiparie, un laboratoire en devint quarante.} \src{Chron. Fabr.} --- \dv{fissipariste} \textit{n.} fondateur issu d'un essaimage.

%EN flammigerist
\E{flammigériste}{fla.mi.ʒe.ʁist}{n.}{lat. \la{flammiger} porte-flamme} Qui porte la flamme d'un projet après le départ de son fondateur ; mainteneur de code ouvert. \ex{Gloire aux fondateurs ; du pain aux flammigéristes.} --- \dv{flammigère} \textit{adj.}

%EN foom
\E*{foom}{fum}{n. m. et v. intr.}{onomatopée angl.} \lbl{vulg.} Auto-amélioration soudaine et explosive d'une intelligence artificielle ; \textit{faire foom}, la subir. Cf. \xr{autocatalyse}, \xr{henotachy}.

%EN fornacist
\E{fornaciste}{fɔʁ.na.sist}{n.}{lat. \la{fornax} fournaise} Gardien des fournaises ; qui veille sur les réacteurs, les fonderies ou les fours, et qu'on honore au calendrier accélérationniste à côté du \xr{pyrogenitor}.

%EN frenocracy
\E{frénocratie}{fʁe.nɔ.kʁa.si}{n. f.}{lat. \la{frenum} bride + gr. κράτος pouvoir} Gouvernement par la bride ; ordre réglementaire dont l'objet est de retenir plutôt que de diriger. --- \dv{frénateur} \textit{n. m.} qui bride ; \dv{frénation} \textit{n. f.}

%EN frigorism
\E{frigorisme}{fʁi.ɡɔ.ʁism}{n. m.}{lat. \la{frigus, frigoris} froid} \lbl{dérog.} Politique de refroidissement délibéré de l'économie ; hiver volontaire du décroissant. \ex{Le frigorisme est une mort thermique prise par anticipation.} \src{Brev. Therm.} --- \dv{frigoriste} \textit{n.} Cf. \xr{meiosist}.

%EN fugacist
\E{fugaciste}{fy.ɡa.sist}{n.}{lat. \la{fugax} fugitif} Qui tient qu'aucun avantage concurrentiel ne dure ; que toute douve est fugitive et doit être recreusée chaque année. \ex{Le fugaciste dort d'un œil, et bâtit chaque jour.} --- \dv{fugacisme} \textit{n. m.}

%EN fulcrarch
\E{fulcrarque}{fyl.kʁaʁk}{n.}{lat. \la{fulcrum} montant de lit, appui + gr. ἀρχός chef} Celui qui règne par le levier ; petite équipe qui meut un grand monde au moyen d'un outil. \ex{Donnez au fulcrarque un modèle et un point d'appui.} \src{Ep. Fabr.} v. --- \dv{fulcrarchie} \textit{n. f.}

%EN fulgurism
\E{fulgurisme}{fyl.ɡy.ʁism}{n. m.}{lat. \la{fulgur} éclair} Doctrine selon laquelle les grands progrès arrivent soudainement, comme la foudre, et que les institutions bâties pour le changement graduel en seront frappées à l'improviste.

%EN fulminatrice
\E{fulminatrice}{fyl.mi.na.tʁis}{n. f.}{fr., formé sur lat. \la{fulminare} lancer la foudre} Machine qui frappe soudain et change tout ce qu'elle touche ; invention de rupture. Cf. \xr{fulgurism}.

%EN funambulocracy
\E{funambulocratie}{fy.nɑ̃.by.lɔ.kʁa.si}{n. f.}{lat. \la{funambulus} danseur de corde + gr. κράτος pouvoir} Gouvernement conçu comme une marche sur la corde raide entre l'accélération et la retenue ; posture du régulateur qui souhaite être loué des deux partis. --- \dv{funambulocrate} \textit{n.}

%EN futurarch
\E{futurarque}{fy.ty.ʁaʁk}{n.}{lat. \la{futurum} l'avenir + gr. ἀρχός chef} Qui prétend gouverner au nom des enfants à naître ; titre disputé entre l'e/acc, qui voudrait bâtir pour eux, et les décels, qui voudraient les protéger. --- \dv{futurarchie} \textit{n. f.}

%EN futurescence
\E{futurescence}{fy.ty.ʁɛ.sɑ̃s}{n. f.}{lat. \la{futurus} qui doit être + -escence, d'après \textit{adolescence}} État de devenir l'avenir ; jeunesse gauche d'une technique entre la démonstration et l'ubiquité. \ex{La machine était dans sa futurescence, maladroite et immortelle.} --- \dv{futurescent} \textit{adj.}

%EN futurophage
\E{futurophage}{fy.ty.ʁɔ.faʒ}{n. et adj.}{lat. \la{futurum} + gr. -φάγος qui mange} \lbl{dérog.} Qui consomme l'avenir pour payer le présent ; épithète que chaque parti applique à l'autre, les accélérationnistes visant ceux qui renoncent à la croissance, les décels ceux qui courent le risque.

%EN gaiaphagy
\E{gaïaphagie}{ɡa.ja.fa.ʒi}{n. f.}{gr. Γαῖα terre + -φαγία action de manger} \lbl{dérog.} Dévoration de la terre ; grief du décel contre l'accélération, à quoi l'accélérationniste répond que la terre est petite et le ciel grand. Cf. \xr{gaiastasis}.

%EN gaiastasis
\E{gaïastase}{ɡa.jas.taz}{n. f.}{gr. Γαῖα terre + στάσις immobilité} Idéal d'une planète tenue à perpétuité dans son état présent, ni améliorée ni exploitée ; paradis du décroissant. --- \dv{gaïastatique} \textit{adj.}

%EN gaiergy
\E{gaïergie}{ɡa.jɛʁ.ʒi}{n. f.}{gr. Γαῖα terre + ἔργον ouvrage} Énergie tirée de la chaleur intérieure de la planète ; la géothermie, considérée comme un don retenu jusqu'à ce que les forages fussent assez profonds. --- \dv{gaïergique} \textit{adj.}

%EN galactophage
\E{galactophage}{ɡa.lak.tɔ.faʒ}{n. et adj.}{gr. γάλα, γάλακτος lait (d'où \textit{galaxie}) + -φάγος qui mange} Mangeur de galaxie ; civilisation du rang \xr{tritobathmic}, conçue comme buvant la Voie lactée. --- \dv{galactophagie} \textit{n. f.}

%EN geitonophobia
\E{géitonophobie}{ʒe.i.tɔ.nɔ.fɔ.bi}{n. f.}{gr. γείτων voisin + φόβος crainte} Crainte des nouveaux voisins, qu'il s'agisse de maisons, de tours, de lignes électriques ou de personnes ; passion animatrice de l'\xr{oikoclast}. \ex{La géitonophobie a bâti plus de murs que la guerre.} \src{Sat. Decel.} xi. --- \dv{géitonophobe} \textit{n. et adj.}

%EN gelotocracy
\E{gélotocratie}{ʒe.lɔ.tɔ.kʁa.si}{n. f.}{gr. γέλως rire + κράτος pouvoir} Gouvernement par le rire ; pouvoir qu'a le mème de rendre une politique ridicule avant qu'elle ne devienne loi. --- \dv{gélotocrate} \textit{n.} Cf. \xr{gelotopoeia}.

%EN gelotopoeia
\E{gélotopée}{ʒe.lɔ.tɔ.pe}{n. f.}{gr. γελωτοποιία action de faire rire} Art de faire des mèmes ; tenu sur l'Agora pour la forme la plus efficace d'argument politique. \ex{Les décels écrivaient des livres blancs ; les accélérationnistes répondirent par la gélotopée, et l'emportèrent.} \src{Agora} --- \dv{gélotopoète} \textit{n.} faiseur de mèmes.

%EN genearchy
\E{généarchie}{ʒe.ne.aʁ.ʃi}{n. f.}{gr. γενεά génération + ἀρχή commandement} Gouvernement par la génération montante ; état d'une société où l'on permet aux jeunes de bâtir avant de leur permettre de voter. Contr. \xr{gerontostasis}.

%EN genesiurge
\E{génésiurge}{ʒe.ne.zjyʁʒ}{n.}{gr. γένεσις origine + -ουργός artisan} Faiseur de commencements ; qui fonde une discipline nouvelle plutôt qu'il n'en avance une ancienne. --- \dv{génésiurgie} \textit{n. f.} Cf. \xr{neoctist}.

%EN genomorphosis
\E{génomorphose}{ʒe.nɔ.mɔʁ.foz}{n. f.}{gr. γένος race, espèce + μόρφωσις façonnement} Édition du génome ; refaçonnement de l'hérédité vivante par un art délibéré. \ex{Après la génomorphose, l'hérédité devint un brouillon.} --- \dv{génomorphique} \textit{adj.} Cf. \xr{biotachy}.

%EN geoperibole
\E{géopéribole}{ʒe.ɔ.pe.ʁi.bɔl}{n. f.}{gr. γῆ terre + περιβολή enceinte} Enceinte d'une planète par des ouvrages orbitaux, anneaux et treillis de satellites ; anticipation protobathmique de l'\xr{helioperibole}.

%EN geophygy
\E{géophygie}{ʒe.ɔ.fi.ʒi}{n. f.}{gr. γῆ terre + φυγή fuite} Fuite hors de la Terre ; départ des premiers colons permanents. \sn{1} L'événement. \sn{2} \lbl{fig.} Tout pas irréversible vers une frontière. --- \dv{géophygique} \textit{adj.}

%EN gerontostasis
\E{gérontostase}{ʒe.ʁɔ̃.tɔs.taz}{n. f.}{gr. γέρων vieillard + στάσις immobilité} Stagnation d'une société vieillissante, où les vieux l'emportent aux voix sur les jeunes et où le titulaire survit au fondateur. --- \dv{gérontostatique} \textit{adj.} Cf. \xr{catastasis}.

%EN gigaergic
\E{gigaergique}{ʒi.ɡa.ɛʁ.ʒik}{adj.}{gr. γίγας géant + ἔργον ouvrage} Se dit des projets qui requièrent des énergies de l'ordre du gigawatt ou au-delà ; des personnes qui pensent à cette échelle. \ex{une ambition gigaergique, un budget mégaergique}.

%EN glossotachy
\E{glossotachie}{ɡlɔ.sɔ.ta.ki}{n. f.}{gr. γλῶσσα langue + τάχος vitesse} Vitesse du changement linguistique ; rythme auquel un mouvement engendre des mots nouveaux, dont ce Lexique est à la fois le registre et le symptôme. --- \dv{glossotachique} \textit{adj.} Cf. \xr{tachylect}.

%EN glossurgy
\E{glossurgie}{ɡlɔ.syʁ.ʒi}{n. f.}{gr. γλῶσσα langue, langage + -ουργία travail} Art d'instruire les machines en langue naturelle ; l'art de l'invite, considéré comme une branche de la rhétorique. \ex{Le glossurge parle, et la machine bâtit.} --- \dv{glossurge} \textit{n.}

%EN gnomarch
\E{gnomarque}{ɡnɔ.maʁk}{n.}{gr. γνώμη jugement, opinion + ἀρχός chef} \lbl{iron.} Qui règne par l'opinion ; chroniqueur dont l'autorité repose sur la prédiction, non sur la production.

%EN gnomonist
\E{gnomoniste}{ɡnɔ.mɔ.nist}{n.}{gr. γνώμων aiguille du cadran solaire} Qui regarde l'ombre tourner ; observateur du progrès qui mesure beaucoup et bâtit peu. Cf. \xr{clepsydromania}.
